\section{Appendix: Scriptural Date and Location}
% Generated. Do not edit.
% Deterministic projection of research/chronology.toml.
% schema: 2
% document: liturgy/roman-rite/postconciliar/roman-missal-third-edition-en-us-2011/propers/temporal/pc-s53-twenty-seventh-sunday-in-ordinary-time-year-a
% generated-by: tools/tpt proper-chronology annotations
\newcommand{\chronologyannotationclaim}[6]{#6}
\newcommand{\chronologyannotationreach}[3]{}
\newcommand{\chronologyannotationgroup}[3]{#3}
\newcommand{\chronologyannotation}[1]{%
  \ifcsname triptychchronologyannotation@#1\endcsname
    \csname triptychchronologyannotation@#1\endcsname
  \else
    \PackageError{triptych}{No chronology annotation for #1}{}%
  \fi}
% appointment-dependency: guidance/liturgy/postconciliar-propers-registry.md sha256=ffe8a6d8b7d79c652bdc49b3d5f6d4dbb32e3506abc7ab0310e745c2b7306ad4
% appointment-dependency: src/gpt/liturgy/roman-rite/postconciliar/roman-missal-third-edition-en-us-2011/propers/registry/README.md sha256=e61f5323480cbe19b7bf3b9ee046cfe8966f93662367513dea9cf9b0f27d5b66
% appointment-dependency: src/gpt/liturgy/roman-rite/postconciliar/roman-missal-third-edition-en-us-2011/propers/registry/formula-dispositions.md sha256=89d9d58ac5bcc07c00d12da3660e07d5ffad72866cc43f40b9581918ed71bde5
% appointment-dependency: src/gpt/liturgy/roman-rite/postconciliar/roman-missal-third-edition-en-us-2011/propers/temporal/pc-s53-twenty-seventh-sunday-in-ordinary-time-year-a/instance/manifest.md sha256=9ab2f377d6de66eb639a2115620d32379b1d22b40cea1ac67dcf429b821630e9
% appointment-dependency: src/gpt/liturgy/roman-rite/postconciliar/roman-missal-third-edition-en-us-2011/propers/temporal/pc-s53-twenty-seventh-sunday-in-ordinary-time-year-a/propers/verified.md sha256=c4367c88e48ca6954481a904edde9c857128f2a94e74649cc38e993cafb22bb8
% appointment-dependency: src/gpt/liturgy/roman-rite/postconciliar/roman-missal-third-edition-en-us-2011/propers/temporal/pc-s53-twenty-seventh-sunday-in-ordinary-time-year-a/research/chronology-inputs.toml sha256=9fd487f1561888df70c71ff005090c7a9ff3a09bf7d846ab70671ac156997841
% appointment-dependency: src/gpt/liturgy/roman-rite/postconciliar/roman-missal-third-edition-en-us-2011/propers/temporal/shared/ordinary-time/weeks/27/propers/verified.md sha256=49ce6830cc5bce8d654908675720c69c73ad0e44ee3832e3196288406a137cd6
% appointment-note: entrance: Esther 13:9-11; identified-basis; vulgate numbering; canonical verse queries, not appointed wording.
% appointment-note: collect: Composed Missal oration; no directly appointed biblical citation.
% appointment-note: first-reading: Isaiah 5:1-7; appointed; vulgate numbering; canonical verse queries, not appointed wording.
% appointment-note: responsorial-psalm: Psalm 80:9, 12, 13-14, 15-16, 19-20; appointed; hebrew numbering; canonical verse queries, not appointed wording.
% appointment-note: responsorial-psalm: Isaiah 5:7a; response; vulgate numbering; whole-verse query envelope, not subverse chronology or appointed wording.
% appointment-note: second-reading: Philippians 4:6-9; appointed; vulgate numbering; canonical verse queries, not appointed wording.
% appointment-note: acclamation: Cf. John 15:16; adaptation; vulgate numbering; canonical verse queries, not appointed wording.
% appointment-note: gospel: Matthew 21:33-43; appointed; vulgate numbering; canonical verse queries, not appointed wording.
% appointment-note: offerings: Composed Missal oration; no directly appointed biblical citation.
% appointment-note: communion-lamentations: Lamentations 3:25; appointed; vulgate numbering; canonical verse queries, not appointed wording.
% appointment-note: communion-corinthians: Cf. 1 Corinthians 10:17; adaptation; vulgate numbering; canonical verse queries, not appointed wording.
% appointment-note: communion-corinthians: 1 Corinthians 10:16; identified-basis; vulgate numbering; canonical verse queries, not appointed wording.
% appointment-note: after-communion: Composed Missal oration; no directly appointed biblical citation.
% entrance: Entrance antiphon
\expandafter\def\csname triptychchronologyannotation@entrance\endcsname{%
  \chronologyannotationgroup{composition}{preferred}{\textbf{Composition}: \chronologyannotationreach{Esth.13.10}{Esth}{true}\chronologyannotationreach{Esth.13.11}{Esth}{true}\chronologyannotationreach{Esth.13.9}{Esth}{true}\chronologyannotationclaim{composition.book-of-esther}{composition}{catholic-traditional-v1}{preferred}{at the end of the reign of Xerxes I (485-465 B.C.) or during the reign of his son Artaxerxes I (465-425 B.C.)}{B.C. 465--425}.}%
}
% first-reading: First reading
\expandafter\def\csname triptychchronologyannotation@first-reading\endcsname{%
  \chronologyannotationgroup{traditional-attribution}{preferred}{\textbf{Traditional attribution}: \chronologyannotationreach{Is.5.1}{Is}{true}\chronologyannotationreach{Is.5.2}{Is}{true}\chronologyannotationreach{Is.5.3}{Is}{true}\chronologyannotationreach{Is.5.4}{Is}{true}\chronologyannotationreach{Is.5.5}{Is}{true}\chronologyannotationreach{Is.5.6}{Is}{true}\chronologyannotationreach{Is.5.7}{Is}{true}\chronologyannotationclaim{israel.prophets.isaias-traditional-era}{traditional-attribution}{catholic-traditional-v1}{preferred}{from the closing year of Ozias, King of Juda (740 B.C.); no prophecies appear to be later than 701 B.C.}{Isaias (ministry in Souvay's traditional account), B.C. 740--701}.}%
}
% responsorial-psalm: Responsorial psalm
\expandafter\def\csname triptychchronologyannotation@responsorial-psalm\endcsname{%
  \chronologyannotationgroup{composition}{nonuniform}{\textbf{Composition} -- No single date applies across every cited locus.}%
}
% second-reading: Second reading
\expandafter\def\csname triptychchronologyannotation@second-reading\endcsname{%
  \chronologyannotationgroup{composition}{disputed}{\textbf{Composition} -- disputed: \chronologyannotationreach{Phil.4.6}{Phil}{true}\chronologyannotationreach{Phil.4.7}{Phil}{true}\chronologyannotationreach{Phil.4.8}{Phil}{true}\chronologyannotationreach{Phil.4.9}{Phil}{true}\chronologyannotationclaim{composition.epistle-to-the-philippians}{composition}{catholic-traditional-v1}{disputed}{(Philemon; Colossians; Ephesians; Philippians), 61}{A.D. 61}; \chronologyannotationreach{Phil.4.6}{Phil}{true}\chronologyannotationreach{Phil.4.7}{Phil}{true}\chronologyannotationreach{Phil.4.8}{Phil}{true}\chronologyannotationreach{Phil.4.9}{Phil}{true}\chronologyannotationclaim{composition.epistle-to-the-philippians}{composition}{catholic-traditional-v1}{disputed}{at Rome (A.D. 62-64)}{A.D. 62--64}.}%
}
% acclamation: Gospel acclamation
\expandafter\def\csname triptychchronologyannotation@acclamation\endcsname{%
  \chronologyannotationgroup{narrated-event}{disputed}{\textbf{Event} -- disputed: \chronologyannotationreach{John.15.16}{Matt.26.20-29 Mark.14.17-25 Luke.22.14-38 John.13.1- John.14 John.15 John.16 John.17.1-26}{true}\chronologyannotationclaim{life-of-christ.last-supper}{narrated-event}{catholic-traditional-v1}{disputed}{the fourteenth day of Nisan}{01-14}.}%
  \space \chronologyannotationgroup{composition}{disputed}{\textbf{Composition} -- disputed: \chronologyannotationreach{John.15.16}{John}{true}\chronologyannotationclaim{composition.gospel-of-john}{composition}{catholic-traditional-v1}{disputed}{from the year 90 to 100 (approximately)}{c. A.D. 90--100}; \chronologyannotationreach{John.15.16}{John}{true}\chronologyannotationclaim{composition.gospel-of-john}{composition}{catholic-traditional-v1}{disputed}{96 or one of the succeeding years}{A.D. 96--100}.}%
}
% gospel: Gospel
\expandafter\def\csname triptychchronologyannotation@gospel\endcsname{%
  \chronologyannotationgroup{narrated-event}{research-pending}{\textbf{Event} -- No narrated-event date in the chronology corpus.}%
  \space \chronologyannotationgroup{composition}{disputed}{\textbf{Composition} -- disputed: \chronologyannotationreach{Matt.21.33}{Matt}{true}\chronologyannotationreach{Matt.21.34}{Matt}{true}\chronologyannotationreach{Matt.21.35}{Matt}{true}\chronologyannotationreach{Matt.21.36}{Matt}{true}\chronologyannotationreach{Matt.21.37}{Matt}{true}\chronologyannotationreach{Matt.21.38}{Matt}{true}\chronologyannotationreach{Matt.21.39}{Matt}{true}\chronologyannotationreach{Matt.21.40}{Matt}{true}\chronologyannotationreach{Matt.21.41}{Matt}{true}\chronologyannotationreach{Matt.21.42}{Matt}{true}\chronologyannotationreach{Matt.21.43}{Matt}{true}\chronologyannotationclaim{composition.gospel-of-matthew}{composition}{catholic-traditional-v1}{disputed}{about A.D. 38-45}{c. A.D. 38--45}; \chronologyannotationreach{Matt.21.33}{Matt}{true}\chronologyannotationreach{Matt.21.34}{Matt}{true}\chronologyannotationreach{Matt.21.35}{Matt}{true}\chronologyannotationreach{Matt.21.36}{Matt}{true}\chronologyannotationreach{Matt.21.37}{Matt}{true}\chronologyannotationreach{Matt.21.38}{Matt}{true}\chronologyannotationreach{Matt.21.39}{Matt}{true}\chronologyannotationreach{Matt.21.40}{Matt}{true}\chronologyannotationreach{Matt.21.41}{Matt}{true}\chronologyannotationreach{Matt.21.42}{Matt}{true}\chronologyannotationreach{Matt.21.43}{Matt}{true}\chronologyannotationclaim{composition.gospel-of-matthew}{composition}{catholic-traditional-v1}{disputed}{about the year 40-42}{c. A.D. 40--42}; \chronologyannotationreach{Matt.21.33}{Matt}{true}\chronologyannotationreach{Matt.21.34}{Matt}{true}\chronologyannotationreach{Matt.21.35}{Matt}{true}\chronologyannotationreach{Matt.21.36}{Matt}{true}\chronologyannotationreach{Matt.21.37}{Matt}{true}\chronologyannotationreach{Matt.21.38}{Matt}{true}\chronologyannotationreach{Matt.21.39}{Matt}{true}\chronologyannotationreach{Matt.21.40}{Matt}{true}\chronologyannotationreach{Matt.21.41}{Matt}{true}\chronologyannotationreach{Matt.21.42}{Matt}{true}\chronologyannotationreach{Matt.21.43}{Matt}{true}\chronologyannotationclaim{composition.gospel-of-matthew}{composition}{catholic-traditional-v1}{disputed}{the years 40-45}{A.D. 40--45}; \chronologyannotationreach{Matt.21.33}{Matt}{true}\chronologyannotationreach{Matt.21.34}{Matt}{true}\chronologyannotationreach{Matt.21.35}{Matt}{true}\chronologyannotationreach{Matt.21.36}{Matt}{true}\chronologyannotationreach{Matt.21.37}{Matt}{true}\chronologyannotationreach{Matt.21.38}{Matt}{true}\chronologyannotationreach{Matt.21.39}{Matt}{true}\chronologyannotationreach{Matt.21.40}{Matt}{true}\chronologyannotationreach{Matt.21.41}{Matt}{true}\chronologyannotationreach{Matt.21.42}{Matt}{true}\chronologyannotationreach{Matt.21.43}{Matt}{true}\chronologyannotationclaim{composition.gospel-of-matthew}{composition}{catholic-traditional-v1}{disputed}{about the year 60-68}{c. A.D. 60--68}; \chronologyannotationreach{Matt.21.33}{Matt}{true}\chronologyannotationreach{Matt.21.34}{Matt}{true}\chronologyannotationreach{Matt.21.35}{Matt}{true}\chronologyannotationreach{Matt.21.36}{Matt}{true}\chronologyannotationreach{Matt.21.37}{Matt}{true}\chronologyannotationreach{Matt.21.38}{Matt}{true}\chronologyannotationreach{Matt.21.39}{Matt}{true}\chronologyannotationreach{Matt.21.40}{Matt}{true}\chronologyannotationreach{Matt.21.41}{Matt}{true}\chronologyannotationreach{Matt.21.42}{Matt}{true}\chronologyannotationreach{Matt.21.43}{Matt}{true}\chronologyannotationclaim{composition.gospel-of-matthew}{composition}{catholic-traditional-v1}{disputed}{about the years 64-67}{c. A.D. 64--67}; \chronologyannotationreach{Matt.21.33}{Matt}{true}\chronologyannotationreach{Matt.21.34}{Matt}{true}\chronologyannotationreach{Matt.21.35}{Matt}{true}\chronologyannotationreach{Matt.21.36}{Matt}{true}\chronologyannotationreach{Matt.21.37}{Matt}{true}\chronologyannotationreach{Matt.21.38}{Matt}{true}\chronologyannotationreach{Matt.21.39}{Matt}{true}\chronologyannotationreach{Matt.21.40}{Matt}{true}\chronologyannotationreach{Matt.21.41}{Matt}{true}\chronologyannotationreach{Matt.21.42}{Matt}{true}\chronologyannotationreach{Matt.21.43}{Matt}{true}\chronologyannotationclaim{composition.gospel-of-matthew}{composition}{catholic-traditional-v1}{disputed}{about the year 50}{c. A.D. 50}.}%
}
% communion-lamentations: Communion antiphon: Lamentations option
\expandafter\def\csname triptychchronologyannotation@communion-lamentations\endcsname{%
  \chronologyannotationgroup{composition}{preferred}{\textbf{Composition}: \chronologyannotationreach{Lam.3.25}{Lam}{true}\chronologyannotationclaim{critical.lamentations.after-jerusalem-destruction}{composition}{catholic-critical-v1}{preferred}{after Jerusalem's destruction in 587 B.C.; grouped later at an unknown date}{Responses to Jerusalem's destruction in 587 B.C.}}%
}
% communion-corinthians: Communion antiphon: Corinthians option
\expandafter\def\csname triptychchronologyannotation@communion-corinthians\endcsname{%
  \chronologyannotationgroup{composition}{disputed}{\textbf{Composition} -- disputed: \chronologyannotationreach{1Cor.10.16}{1Cor}{true}\chronologyannotationreach{1Cor.10.17}{1Cor}{true}\chronologyannotationclaim{composition.first-epistle-to-the-corinthians}{composition}{catholic-traditional-v1}{disputed}{(1 and 2 Corinthians; Galatians), 56}{A.D. 56}; \chronologyannotationreach{1Cor.10.16}{1Cor}{true}\chronologyannotationreach{1Cor.10.17}{1Cor}{true}\chronologyannotationclaim{composition.first-epistle-to-the-corinthians}{composition}{catholic-traditional-v1}{disputed}{about Easter A.D. 57}{A.D. 57}.}%
}

The dates below distinguish composition, traditional attribution and narrated event. The Catholic Encyclopedia proposals are dated historical judgments, not a present scholarly consensus. Book-level dates do not establish the age of a liturgical adaptation or every constituent passage. The entries follow canonical order; Isaiah's response is consolidated with its first-reading locus. The composed orations have no biblical dates.

\begin{dossiertable}
Entrance & Esther 13:9--11; Missal cf. 4:17 & Persian court narrative at Susan; exact prayer setting not independently dated. & \chronodate{entrance}{\chronologyannotation{entrance}}\\
\dossierprose{McMahon's traditional book-level proposal is concessive and calls authorship conjectural. It does not date the exact wording of Mordecai's addition. The court narrative's location and the book's composition remain different questions.}
\midrule
Psalm & Psalm 80 (Vulgate 79), selected verses; response Isaiah 5:7a & Israel recalled from Egypt; sea and river in the poem. & \chronodate{responsorial-psalm}{\chronologyannotation{responsorial-psalm}}\\
\dossierprose{The psalm's broad pre-Maccabean composition limit in the NABRE introduction does not securely date this individual psalm; the response has Isaiah's different attribution horizon. Augustine's Jordan and Bellarmine's Euphrates are distinct interpretations, not an agreed composition location.}
\midrule
First reading and response & Isaiah 5:1--7; 5:7a & Jerusalem and Judah addressed; Israel named as vineyard. & \chronodate{first-reading}{\chronologyannotation{first-reading}}\\
\dossierprose{Souvay dates Isaiah's ministry in the traditional account. This is neither the composition date of chapter 5 nor the date of a particular performance of the song. His endpoint is the apparent last prophecy, not a known date of death.}
\midrule
Communion alternative & Lamentations 3:25 & Jerusalem's ruin is the communal horizon; the speaker's precise identity is not established. & \chronodate{communion-lamentations}{\chronologyannotation{communion-lamentations}}\\
\dossierprose{The critical fallback in the NABRE introduction gives a lower composition boundary after the destruction, not a date for this verse. Traditional Jeremiah attribution differs from that introduction's judgment that the poems are anonymous. Separate poems may have been grouped later at an unknown time.}
\midrule
Gospel & Matthew 21:33--43 & Jesus' controversy in Jerusalem's temple (21:23); composition setting distinct. & \chronodate{gospel}{\chronologyannotation{gospel}}\\
\dossierprose{Jacquier's proposals depend upon differing testimonies: the early dispersal premise is unreliable, the Eusebian reckoning conditional, and the Irenaean evidence inconclusive. Durand's proposal concerns an Aramaic original and leaves the Greek rendering undated. These are not six secure dates for the same Greek text.}
\midrule
Acclamation & Cf. John 15:16 & Farewell discourse to the disciples; proposed Gospel composition at Ephesus is a different location. & \chronodate{acclamation}{\chronologyannotation{acclamation}}\\
\dossierevent{Narrated event: 01-14 means 14 Nisan, not 14 January; no year is supplied. Maas's disputed Supper reconstruction binds John 13--17 broadly and does not locate each sentence in an exact room.}
\dossierprose{Durand and Fonck supply separate late-first-century composition proposals. Fonck's Ephesus concerns composition, not the place where Jesus speaks. The event date belongs to Maas's complete Passion reconstruction and cannot be spliced into a different reconstruction.}
\midrule
Communion alternative & Cf. 1 Corinthians 10:17; context 10:16 & Corinth is the audience; Aherne locates Paul's correspondence in his Ephesian mission. & \chronodate{communion-corinthians}{\chronologyannotation{communion-corinthians}}\\
\dossierprose{Prat's Pauline table and Aherne's approximate Easter dating differ. They date the epistle, not the institution of the Eucharist or the antiphon's compilation.}
\midrule
Second reading & Philippians 4:6--9 & Philippi is the audience; the historical sources locate composition in Roman imprisonment. & \chronodate{second-reading}{\chronologyannotation{second-reading}}\\
\dossierprose{Prat and Vander Heeren give different imprisonment chronologies. Vander Heeren's confident language is his 1911 position, not current unanimity. The place of writing is not the recipients' city or an exact prison room.}
\end{dossiertable}
